\chapter{Image Compression --- Solutions}

%% -------------------------------------------------------------------------
\section*{Lay-McDonald \S7.5: Covariance and Principal Components}

\begin{exercise}{7.5.1--2}
Convert the matrix of observations to mean-deviation form and construct the sample covariance matrix.
\end{exercise}
\begin{solution}
\textbf{Ex 1:} $X=\begin{pmatrix}19&22&6&3&2&20\\12&6&9&15&13&5\end{pmatrix}$ ($2\times6$). Sample mean $M=\frac{1}{6}\sum_{k=1}^6 X_k = \begin{pmatrix}12\\10\end{pmatrix}$. Mean-deviation: $B=X-M$ (each column $\hat{X}_k=X_k-M$). Covariance: $S=\frac{1}{N-1}BB^T=\frac{1}{5}BB^T$.
\begin{lstlisting}
X = [19 22 6 3 2 20; 12 6 9 15 13 5];
M = mean(X,2); B = X - M;
S = (1/(size(X,2)-1)) * B*B';
\end{lstlisting}
\textbf{Octave output:} $M=(12,10)^T$, $S=\begin{pmatrix}86&-27\\-27&16\end{pmatrix}$.

\textbf{Ex 2:} $X=\begin{pmatrix}1&5&2&6&7&3\\3&11&6&8&15&11\end{pmatrix}$. $M=(4,9)^T$. \textbf{Octave output:} $S=\begin{pmatrix}5.6&8\\8&18\end{pmatrix}$.
\end{solution}

\begin{exercise}{7.5.3--4}
Find the principal components of the data for Exercises 1 and 2.
\end{exercise}
\begin{solution}
Principal components = unit eigenvectors of $S$ ordered by decreasing eigenvalue. First PC $u_1$ maximizes variance of $y_1=u_1^TX$.
\begin{lstlisting}
[V,D] = eig(S); [d,ord] = sort(diag(D),'descend');
V = V(:,ord); u1 = V(:,1);
\end{lstlisting}
\textbf{Octave output:} Ex 3 (data from Ex 1): eigenvalues $\lambda_1\approx95.20$, $\lambda_2\approx6.80$; $u_1\approx(-0.95, 0.32)^T$. Ex 4: $\lambda_1\approx21.92$, $\lambda_2\approx1.68$; $u_1\approx(0.44, 0.90)^T$.
\end{solution}

\begin{exercise}{7.5.5--6}
Given covariance matrix $S$, find first principal component and \% of total variance.
\end{exercise}
\begin{solution}
First PC = unit eigenvector $u_1$ for largest eigenvalue $\lambda_1$. \% variance $=\lambda_1/\operatorname{tr}(S)$.
\begin{lstlisting}
S5 = [164.12 32.73 81.04; 32.73 539.44 249.13; 81.04 249.13 189.11];
[V5,D5] = eig(S5); [d5,ord5] = sort(diag(D5),'descend');
pct5 = d5(1)/sum(d5);  % Ex 5
S6 = [29.64 18.38 5.00; 18.38 20.82 14.06; 5.00 14.06 29.21];
[V6,D6] = eig(S6); [d6,ord6] = sort(diag(D6),'descend');
pct6 = d6(1)/sum(d6);  % Ex 6
\end{lstlisting}
\textbf{Octave output:} Ex 5 (Landsat Homestead): first PC explains \textbf{75.90\%} of variance. Ex 6 (Columbia River): first PC explains \textbf{64.89\%} of variance.
\end{solution}

%% -------------------------------------------------------------------------
\section*{Justin's Notes: Exercises 10.1--10.6}

\begin{exercise}{10.1}
Consider the $4\times4$ image (shades 0, 0.5, 1). (a) Enter into matrix $A$. (b) Find SVD. (c) Simplify with 3, 2, 1 singular values; draw best representation.
\end{exercise}
\begin{solution}
\textbf{(a)} From the image in the text (black/white/grey pattern), a plausible $4\times4$ matrix using only 0, 0.5, 1:
\[
A = \begin{pmatrix}1&0&1&1\\1&1&0&1\\0&1&0&1\\0&1&1&0\end{pmatrix}
\]
(Adjust entries to match the actual figure if different.)

\textbf{(b)}
\begin{lstlisting}
A = [1 0 1 1; 1 1 0 1; 0 1 0 1; 0 1 1 0];
[U,S,V] = svd(A);
svals = diag(S);
fprintf('Singular values: '); disp(svals');
\end{lstlisting}
\textbf{Octave output:} $\sigma_1\approx2.60$, $\sigma_2\approx1.26$, $\sigma_3\approx1.18$, $\sigma_4\approx0.52$. For $k=3,2,1$: variance preserved 97.34\%, 83.35\%, 67.40\%.

\textbf{(c)} For $k=3,2,1$: set $S(i,i)=0$ for $i>k$, then $A'=U*S*V'$. Clamp values to $[0,1]$ for display (or leave as-is per notes). Draw the resulting $4\times4$ grid: darker = lower value, lighter = higher.
\begin{lstlisting}
for k = [3 2 1]
  Sp = S; Sp((k+1):4, (k+1):4) = 0;
  Ap = U*Sp*V';
  fprintf('k=%d, variance preserved: %.4f\n', k, ...
    sum(svals(1:k).^2)/sum(svals.^2));
  disp(Ap);
end
\end{lstlisting}
\end{solution}

\begin{exercise}{10.2}
Consider the $5\times5$ image (shades 0, 0.5, 1). (a) Enter into matrix $A$. (b) Find SVD. (c) Simplify with 4, 3, 2, 1 singular values; draw.
\end{exercise}
\begin{solution}
\textbf{(a)} From the image in the text:
\[
A = \begin{pmatrix}
1&0.5&0&1&1\\
0.5&1&0.5&0&1\\
0&0.5&1&0.5&0\\
1&0&0.5&1&0.5\\
1&1&0&0.5&1
\end{pmatrix}
\]
(Adjust to match the actual figure.)

\textbf{(b)--(c)} Same procedure. \textbf{Octave output:} $\sigma_1\approx3.16$, $\sigma_2\approx1.21$, $\sigma_3\approx1.17$, $\sigma_4\approx0.34$, $\sigma_5\approx0.21$. Variance preserved: $k=4$: 99.67\%; $k=3$: 98.80\%; $k=2$: 88.22\%; $k=1$: 77.01\%.
\begin{lstlisting}
A = [1 0.5 0 1 1; 0.5 1 0.5 0 1; 0 0.5 1 0.5 0; 1 0 0.5 1 0.5; 1 1 0 0.5 1];
[U,S,V] = svd(A);
svals = diag(S);
for k = [4 3 2 1]
  Sp = S; for i=k+1:5, Sp(i,i)=0; end
  Ap = U*Sp*V';
  fprintf('k=%d: var=%.4f\n', k, sum(svals(1:k).^2)/sum(svals.^2));
end
\end{lstlisting}
\end{solution}

\begin{exercise}{10.3}
Load a square image ($\sim200\times200$). (a) Load, grayscale, scale to $[0,1]$. (b) SVD. (c) Count singular values. (d)--(f) Preserve 75\%, 50\%, 25\% of singular values; report quality \%. (g) Min $k$ for 99.9\% quality; compression ratio; print result.
\end{exercise}
\begin{solution}
\begin{lstlisting}
A = imread('image.jpg');
A = rgb2gray(A); A = im2double(A);
A = A - min(A(:)); A = A / max(A(:));

[U,S,V] = svd(A);
svals = diag(S);
n = length(svals);
fprintf('(c) Singular values: %d\n', n);

for pct = [0.75 0.50 0.25]
  k = round(pct * n);
  var_pres = sum(svals(1:k).^2) / sum(svals.^2);
  fprintf('Preserve %.0f%% of values: k=%d, quality=%.4f%%\n', ...
    pct*100, k, var_pres*100);
end

target = 0.999;
cum = cumsum(svals.^2);
tot = cum(end);
k_min = find(cum/tot >= target, 1);
var_at_k = cum(k_min)/tot;
comp_ratio = 2*k_min*n / n^2;
fprintf('(g) Min k for 99.9%%: %d, quality=%.4f%%, compression=%.2f\n', ...
  k_min, var_at_k*100, comp_ratio);

Sp = S; Sp((k_min+1):end, (k_min+1):end) = 0;
Ap = U*Sp*V';
imshow(Ap, 'border', 'tight');
\end{lstlisting}
\end{solution}

\begin{exercise}{10.5}
For each 10$\times$10 image with given singular values, find (a)--(d) the minimum $k$ for 99.9\% variance and the compression ratio.
\end{exercise}
\begin{solution}
Quality $= \frac{\sigma_1^2+\cdots+\sigma_k^2}{\sigma_1^2+\cdots+\sigma_{10}^2}$. Find smallest $k$ with $\geq 0.999$. Compression ratio $= 2k/n = 2k/10 = k/5$.

\textbf{(a)} $\sigma = [80.26, 63.35, 20.59, 8.47, 2.88, 1.68, 0.80, 0.69, 0.66, 0.65]$.
\begin{lstlisting}
s = [80.2608 63.3520 20.5871 8.4696 2.8841 1.6763 0.7962 0.6926 0.6553 0.6520];
cum = cumsum(s.^2); k = find(cum/cum(end) >= 0.999, 1);
fprintf('k=%d, ratio=%.2f\n', k, 2*k/10);  % ratio = 2k/n
\end{lstlisting}
\textbf{Octave output:} (a) $k=5$, ratio $=1.00$; (b) $k=4$, ratio $=0.80$; (c) $k=3$, ratio $=0.60$; (d) $k=5$, ratio $=1.00$. (When $k=5$ or more for $n=10$, ratio $\geq 1$ so no compression benefit.)
\end{solution}

\begin{exercise}{10.6}
For an $m\times n$ image simplified with $k$ singular values: (a) compression ratio; (b) criteria for worth doing.
\end{exercise}
\begin{solution}
\textbf{(a) Storage.}

We store:
\begin{itemize}
  \item $k$ left singular vectors $\bar{u}_1,\ldots,\bar{u}_k$, each in $\mathbb{R}^m$: $km$ values.
  \item $k$ right singular vectors $\bar{v}_1,\ldots,\bar{v}_k$, each in $\mathbb{R}^n$: $kn$ values.
  \item $k$ singular values: $k$ values (often negligible for large $m,n$).
\end{itemize}
Total: $km + kn = k(m+n)$ values (ignoring the $k$ singular values for the ratio).

Original image: $mn$ values.

\textbf{Compression ratio:}
\[
\boxed{\frac{k(m+n)}{mn} = k\left(\frac{1}{n}+\frac{1}{m}\right).}
\]

\textbf{(b) Worth doing when} $k(m+n) < mn$, i.e.,
\[
\boxed{k < \frac{mn}{m+n}.}
\]

For square $n\times n$: ratio $= 2k/n$, worth when $k < n/2$. $\square$
\end{solution}
